% !TeX root = ../report.tex
\documentclass[../report.tex]{subfiles}
\begin{document}

\section{串联校正装置的设置和性能分析}
本节第一部分进行串联超前校正装置的设计和分析，第二部分进行迟后，超前校正
装置的设计和分析，第三部分进行串联迟后校正装置的设计和分析。

\subsection{串联超前校正装置的设计和分析}
串联超前校正有如下形式：
\begin{equation}
    G_c(s) = \frac{aTs+1}{Ts+1}
\end{equation}

串联超前校正装置的特点是相角超前，幅值增加，适用于相位裕度和幅值裕度都不足
的系统。此系统相位裕度不足，故而可以使用超前校正。
首先确定需要超前的相角$\phi_m$:
\[
\begin{split}
    \phi_m &= \gamma^*-\gamma_0 +(5\sim 10)^\circ \\
    &= 45^\circ - 17.96^\circ + 7^\circ = 34.04^\circ
\end{split}
\]
$\phi_m< 60^\circ$,可以使用超前校正。继续求解得到
\[
    a = \frac{1+\sin \phi_m}{1 - \sin \phi_m} = 3.543
\]
校正装置的两个转折点的中点对应校正后系统的截止频率$\omega_c$
,由于校正后系统在此处增益为$10\lg a = 5.49$，所以需要找到
原系统幅频为$|G(jw) = -5.49|$的点，也就是$\omega_c = 8.75$
那么，两个转折点$\omega_C = \frac{1}{aT}$和$\omega_D = \frac{1}{T}$
的中点$\omega_D$有：
\[
    \omega_c = \sqrt{\omega_C \omega_D} = \frac{1}{\sqrt{aT^2}}=8.75
\]
得到，$T = 0.0861$
于是，校正装置就被确定为：
\[
    G_c(s) = \frac{\frac{s}{3.278}+1}{\frac{s}{11.6144}+1}
\]
经过计算，此时系统的相角裕度是$GM = 43.14^\circ$。
性能已经接近指标，不妨向右移动第二个转折点，增大校正幅度。
校正装置为：
\[
    G_c(s) = \frac{\frac{s}{3.278}+1}{\frac{s}{13.6144}+1}
\]
此时系统的相角裕度是$GM = 46.64^\circ$。
校正系统的伯德图如图\ref{fig:for}，校正后系统的伯德图如图\ref{fig:for_sys}
，阶跃响应是\ref{fig:for_time}。此时，系统幅值裕度依然是无穷，时域上
调节时间减少（0.49s），而且超调量大幅降低（24.2\%）。
\begin{figure}[H]
    \centering
    \includegraphics[width=0.7\textwidth]{./ques2_2/for.png}
    \caption{校正装置伯德图}
    \label{fig:for}
\end{figure}
\begin{figure}[H]
    \centering
    \includegraphics[width=0.7\textwidth]{./ques2_2/for_sys.png}
    \caption{校正后系统伯德图}
    \label{fig:for_sys}
\end{figure}
\begin{figure}[H]
    \centering
    \includegraphics[width=0.7\textwidth]{./ques2_2/for_time.png}
    \caption{校正后系统阶跃响应}
    \label{fig:for_time}
\end{figure}

\subsection{串联迟后-超前校正的设计和分析}

串联迟后-超前校正具有形式：
\begin{equation}
    G_c(s) = \frac{T_a s +1}{aT_a s +1}\frac{T_b s +1}{\frac{T_b}{a}s +1}
\end{equation}

串联迟后-超前校正的特点是幅值衰减，相角超前，综合利用了迟后网络和超前
网络的特性，提高系统性能。

同样地，先确定$\phi_m$:
\[
    \begin{split}
          \phi_m &= \gamma^* -\gamma_0(\omega_c^*)+6^\circ \\ 
          &= 45^\circ -17.96^\circ +6^\circ = 33.04^\circ
    \end{split}
\]
那么
\[
    a = \frac{1+\sin \phi_m}{1-\sin \phi_m} = 3.3978
\]
类似于超前校正中确定转折点的方法，可以得到$|G(j\omega)| = 10 \lg a = 5.312$
处就是$\omega_C$和$\omega_D$的中点，即$\omega = 4.46323$，于是
\[
    \begin{cases}
    \omega_c = \sqrt{\omega_C \omega_D} = \frac{1}{\sqrt{aT^2}}=4.46323\\
    \omega_c = \frac{1}{aT}
    \omega_d = \frac{1}{T}
    \end{cases}
\]
其中$T = \frac{T_b}{a}$，解得，$\omega_C = 2.4231$，$\omega_D = 8.2271$。
又由于$\omega_E$和$\omega_c$相距十倍频程，
\[
    \begin{cases}
    \omega_E &= 0.1\omega_c\\
    \omega_F &= \frac{\omega_E}{a}
    \end{cases}
\]
解得
$\omega_a = 0.131356$， $\omega_b = 0.446323$

该系统表达式
\[
    G_c(s) =\frac{\frac{s}{0.446323}+1}{\frac{s}{0.131356}+1}\frac{\frac{s}{2.4231}+1}{\frac{s}{8.2271}+1}
\]
校正系统的伯德图如图\ref{fig:bf}，校正后系统的伯德图如图\ref{fig:bf_sys}
，阶跃响应是\ref{fig:bf_time}。此时，系统相角裕度是$53.17^\circ$，幅值裕度依然是无穷，时域上
调节时间减少（2.34s），而且超调量大幅降低（18.7\%）。
\begin{figure}[H]
    \centering
    \includegraphics[width=0.7\textwidth]{./ques2_2/bf.png}
    \caption{校正装置伯德图}
    \label{fig:bf}
\end{figure}
\begin{figure}[H]
    \centering
    \includegraphics[width=0.7\textwidth]{./ques2_2/bf_sys.png}
    \caption{校正后系统伯德图}
    \label{fig:bf_sys}
\end{figure}
\begin{figure}[H]
    \centering
    \includegraphics[width=0.7\textwidth]{./ques2_2/bf_time.png}
    \caption{校正后系统阶跃响应}
    \label{fig:bf_time}
\end{figure}

\subsection{串联迟后校正系统的设计和分析}
串联迟后系统具有形式：
\begin{equation}
   G_c(s) = \frac{bTs+1}{Ts +1}
\end{equation}
串联迟后校正具有相角迟后，幅值衰减的特性，能够挖掘系统的相角储备，适合截止频率有余裕，但相角裕度条件不足且无法通过
超前校正弥补的情况。
原系统中不满足相角条件要求但可以用超前校正弥补，而且截止频率却并不充足（调节时间较长），并不适配于串联迟后校正。但是如果只是满足系统的幅值，相角裕度
的条件，那么串联迟后校正也是可以使用的。

详细步骤如下：
首先选定系统的新截止频率$\omega_c$,
新截止频率需要满足：
\[
45^\circ + 6^\circ = 180^\circ + \angle G(j\omega_c) 
\]
加上$6^\circ$的原因是迟后校正系统第二个转折点十倍频程后会引入一个不超过$6^\circ$的相角。
于是选定 $\angle G(j\omega_c) = 129^\circ$。在图\ref{fig:back1}中寻找此点，为$\omega_c = 1.619$
\begin{figure}[H]
    \centering
    \includegraphics[width=0.8\textwidth]{./ques2_2/back1.png}
    \caption{原系统的相频曲线}
    \label{fig:back1}
\end{figure}

接着确定需要降低的幅值，也就是$20\lg b$，就是图\ref{fig:back2}中对应于$\omega_c$的幅值的相反数。
\begin{figure}[H]
    \centering
    \includegraphics[width=0.8\textwidth]{./ques2_2/back2.png}
    \caption{原系统的幅频曲线}
    \label{fig:back2}
\end{figure}
于是
\[
20\lg b = -19.648 \Rightarrow  b = 0.1041
\]

于是，容易得到两个转折频率$\omega_E$、$\omega_F$：
\[
\begin{split}
    \omega_F &= 0.1 \omega_c\\
    \omega_E &= b\omega_F
\end{split}
\]

得到$\omega_E = 0.0207036$、$\omega_F = 0.2277$。
迟后校正系统为：
\[
    G_c(s) = \frac{\frac{s}{0.2277}+1}{\frac{s}{0.0207036}+1}
\]
系统伯德图如图\ref{fig:back}
\begin{figure}[H]
    \centering
    \includegraphics[width=0.7\textwidth]{./ques2_2/back.png}
    \caption{校正装置伯德图}
    \label{fig:back}
\end{figure}

校正后系统的伯德图如图\ref{fig:back_sys}，阶跃响应如图\ref{fig:back_time}
\begin{figure}[H]
    \centering
    \includegraphics[width=0.7\textwidth]{./ques2_2/back_sys.png}
    \caption{校正后系统伯德图}
    \label{fig:back_sys}
\end{figure}

\begin{figure}[H]
    \centering
    \includegraphics[width=0.7\textwidth]{./ques2_2/back_time.png}
    \caption{校正后系统阶跃响应}
    \label{fig:back_time}
\end{figure}
校正后系统的相角裕度是$45.56^\circ$，幅值裕度不变，满足设计要求。但是从时域响应图\ref{fig:back_time}中
可以看出，由于截止频率的减小，系统响应速度变慢，系统调节时间并没有因为震荡被遏制而得到缩减，反而明显增大（7.6s）；但系统的超调量明显减少（28.2\%）。


\subsection{三种校正方式对比分析}
针对同一被控对象，分别设计了超前校正、滞后校正及滞后-超前校正装置。三种校正策略在改善系统性能时的侧重点各不相同，具体表现在频域指标和时域响应上的差异。三种校正方式校正后系统频域特性如图\ref{fig:bode}，时域响应如图\ref{fig:time}。

\begin{figure}[H]
    \centering
    \includegraphics[width=0.8\textwidth]{./ques2_2/bode.png}
    \caption{伯德图对比}
    \label{fig:bode}
\end{figure}

\begin{figure}[H]
    \centering
    \includegraphics[width=0.8\textwidth]{./ques2_2/time.png}
    \caption{阶跃响应对比}
    \label{fig:time}
\end{figure}




从设计原理上看，\textbf{超前校正}适用于稳态精度已满足而动态性能不足的场合。它利用网络在中频段的相位超前特性，提高相角裕度，减小阶跃响应的超调量，并加快响应速度（增大剪切频率 $\omega_c$）。但会抬升高频段增益，降低对高频噪声的抗干扰能力。

\textbf{滞后校正}适用于动态性能满足而稳态精度不足的场合。它通过将极点配置在低频段，提高开环直流增益以降低稳态误差；同时将剪切频率左移，牺牲响应速度换取相角储备，保证稳定性。

\textbf{滞后-超前校正}适用于对稳态精度和动态性能均有较高要求的场合。该方案结合两者优势：滞后部分提高低频增益改善稳态性能，超前部分补偿中频相位提升动态响应品质。虽然物理实现相对复杂，但在单一校正无法满足全部性能指标时，具有最佳的调节灵活性。

综上所述，三种校正方式的性能特点对比如表~\ref{tab:correction_comparison}所示。

\begin{table}[htbp]
    \centering
    \caption{三种校正方式性能对比}
    \label{tab:correction_comparison}
    \begin{tabular}{lccc}
        \toprule
        \textbf{性能指标} & \textbf{超前校正} & \textbf{滞后校正} & \textbf{滞后-超前校正} \\
        \midrule
        超调量 $\sigma\%$ & 减小 & 略有减小 & 减小 \\
        调节时间 $t_s$ & 缩短 & 延长 & 适中 \\
        剪切频率 $\omega_c$ & 增大（响应快） & 减小（响应慢） & 适中 \\
        抗高频噪声 & 较差 & 较好 & 一般 \\
        实现复杂度 & 简单 & 简单 & 较复杂 \\
        \bottomrule
    \end{tabular}
\end{table}

就本题目而言，由于超前校正对高频增益不大，采用超前校正，不仅实现简单，而且明显减少超调量和调节时间，
迟后-超前校正表现适中，校正装置更加灵活可调，但是实现较为复杂。迟后系统只是勉强达到基本设计要求，
对超调的抑制明显不如前两者，而且调节时间相比前两者更长，不适宜应用在此系统中。故，如果对高频噪声
不敏感，应当使用超前校正；如果不考虑实现成本，可以使用迟后-超前校正。



% \begin{figure}[htbp]
%     \centering
%     % 假设图片文件名为 comparison.png
%     \includegraphics[width=0.8\textwidth]{comparison.png}
%     \caption{三种校正后的系统阶跃响应对比}
%     \label{fig:step_response_comparison}
% \end{figure}

\end{document}

